\begin{helper}
For every fixed real number \(c>0\),
\[
c\,\noderef{RamseyScale}(k)\to\infty
\]
as \(k\to\infty\).
\end{helper}
\begin{proof}
Fix \(c>0\).  By the definition of \(\noderef{RamseyScale}\),
\[
c\,\noderef{RamseyScale}(k)=c\,\frac{k^2}{\log k}.
\]
We prove convergence to \(+\infty\) from the order characterization of
tending to infinity.  Let \(B\in\mathbb R\) be arbitrary.  Choose \(K_1\) so
large that \(k\ge K_1\) implies \(k>1\).  Then \(\log k>0\).  Also, by the
standard logarithmic estimate \(\log x\le x\) for \(x>0\), the same range has
\(\log k\le k\).  Since both denominators are positive,
\[
\frac{k^2}{\log k}\ge \frac{k^2}{k}=k .
\]
If \(B\le0\), then for all \(k\ge K_1\),
\[
c\,\frac{k^2}{\log k}>0\ge B.
\]
If \(B>0\), choose \(K_2\) so large that \(cK_2\ge B\), and set
\(K=\max\{K_1,K_2\}\).  For every \(k\ge K\),
\[
c\,\frac{k^2}{\log k}\ge ck\ge cK_2\ge B.
\]
Thus for every real lower bound \(B\), the values
\(c\,\noderef{RamseyScale}(k)\) are eventually at least \(B\).  This is exactly
\[
c\,\noderef{RamseyScale}(k)\to\infty .
\]
\end{proof}
